A bounded coordinated transaction realizes the previously blocked critical witness. It reserves the target suffix, removes the blocker subtrees serving the ten affected sinks, and retains all paths to unaffected sinks. Blocker networks and sinks are processed in ascending ID/pin order; each attachment uses the current occupancy and recomputed retained-tree prefixes. A failed attachment would discard the transaction. The tree operations pass 1,000 randomized independent path-preservation, ancestry, pruning and attachment checks. All ten attachments succeed, and 91 new directed edges are checked against the original RRG. Unrelated network traces are unchanged and all terminals are restored before native analysis.

The unchanged VPR binary passes native FULL legality and full STA. CPD improves from 6.20393 to 6.09618\,ns (1.736802\%); setup TNS improves from $-35113.1$ to $-35066.8$\,ns. The target RR delay decreases by 0.379291\,ns. Eleven sink delays change: eight decrease and three increase, including blocker degradations of approximately 0.274651, 0.646302 and 0.460476\,ns. Thus final acceptance is based on the whole-design STA result, not target improvement alone. Occupied SB pairs increase from 119 to 128, CB crossing output sources from 956 to 959, and planar wirelength from 228,420 to 228,481. The worst path moves to bram\_select[4] to bram\_wdata\_c\_5\_ext[53]. This is a feasible coordinated improvement on one current state, not a globally optimal or promoted method.

Prototype construction takes 48.047812\,s including preparation and physical audits. Its ten query cores sum to 5.395719\,s; the ten complete query processes sum to 13.894189\,s, including repeated static-graph loading. The native analysis process takes 8.915268\,s. These are nested or distinct stage measurements as stated, not additive core/parent timers, not full initial-route process measurements, and not a formal native speed comparison. Persistent graph reuse and in-process integration remain engineering work rather than claimed speedups.

The next current-STA free-capacity probe is applied to the repaired Tf state and to the matched strong A0/R1 states. Packed netlist, placement and RRG digests match, but the route states and prior repair opportunities differ. The four Tf targets contain one positive witness (net 7647 pin 56, 0.662117\,ns local gain); A0's five targets contain none; R1's five contain one (net 6590 pin 4, 0.185826\,ns local gain). Both positive proposals pass native FULL/STA. Tf improves further from 6.09618 to 6.01928\,ns and is accepted. R1 stays at 6.01810\,ns and is rejected by the registered strict-CPD rule, despite only its target physical delay decreasing. R1's reported setup TNS also stays at $-33483.7$\,ns while the worst-path start moves from bram\_select[7] to bram\_select[6], exposing a tied-path plateau rather than evidence that the physical repair did not happen.

The Tf trajectory now improves 3.041823\% from its frozen 6.20812\,ns state but remains 0.00118\,ns worse than the strong A0/R1 CPD. Its final SB/CB counts are 129/960 and planar wirelength is 228,483; method superiority is not established. This stage performs 24 new queries and three legal native analysis calls, with two accepted interventions and one legal strict-CPD rejection. The following separate acceptance-factor experiment tests: permit a zero-CPD-gain single-connection step only when all audited sink physical delays are nonincreasing with at least one strict decrease and native setup/hold metrics do not worsen, then test whether subsequent steps reduce CPD. Existing strict-policy evidence remains immutable. This zero-RC componentwise condition does not automatically apply to coordinated trades, nonzero-RC loading, or broader timing objectives. Equal-opportunity complete-policy A0/R1/Tf comparisons are still required.
